\section{Empirical evaluation}
\label{sec:minimum-results}
Ten UCI classification tasks with 4--41 predictors supplied 30 outer folds.
Six wider tasks with 72--856 predictors supplied 16, including chronological
splits for Ozone and SECOM, participant-level splits for Parkinson and the
official Urban split. Development splits allocated 75\%, 12.5\% and 12.5\%
to fitting, nomination and the refit check, by distinct dates in chronological
order for Ozone and SECOM and by stratified rows elsewhere. We ranked split-used features by XGBoost gain
and took the shortest tested prefix whose masked validation scores met the
8\% relative allowance as the anchor. We kept its leading half as one
candidate and sampled up to seven other subsets of the same size, rounding
the retained width up. Eleven singleton-anchor folds offered no deletion and remained in all panel means.

We report the fraction of anchor columns removed by the delivered model,
counting zero when it breaches an outer tolerance or returns the anchor.
We denote this effective compression by \(U\) and report the breach rate
separately. We averaged folds within tasks and gave each task equal weight.
Paired 95\% intervals resample whole tasks 30,000 times with their fitted
models and folds held fixed.

\begin{table}[tb]
\centering\small
\caption{Nomination cost and outer performance on identical Gain candidate
sets. Random forests performed both masking and refitting. \(U\) and breach
are task-weighted percentages. Masks count distinct predictions on the
nomination sample. Seconds sum within-fold medians from three repetitions.
The Gain prefix and uniform choice need no masked prediction.}
\label{tab:minimum-cost-outcome}
\setlength{\tabcolsep}{4pt}
\begin{tabular}{@{}lrrrrrrrr@{}}
\toprule
&\multicolumn{4}{c}{Ten tasks, 30 folds}&\multicolumn{4}{c}{Six wider tasks, 16 folds}\\
\cmidrule(lr){2-5}\cmidrule(l){6-9}
Rule & Masks & Sec. & \(U\) & Breach & Masks & Sec. & \(U\) & Breach\\
\midrule
Gain prefix     &0   &--- &1.56&3.33 &0     &--- &2.72 &16.67\\
Uniform choice  &0   &--- &4.79&6.81 &0     &--- &8.21 &12.96\\
Single point     &126 &0.61&6.33&10.00&125 &0.65&13.77&16.67\\
Depth-first vote &189 &0.89&2.89&3.33 &366 &1.96&8.28 &16.67\\
Best-first vote  &170 &0.81&2.89&3.33 &300 &1.63&8.28 &16.67\\
\bottomrule
\end{tabular}
\end{table}

\paragraph{Nomination cost.}
Depth-first pruning visits candidates in tie order and abandons each row as soon
as its upper bound cannot beat the incumbent. We kept both orders fixed, fitted a 100-tree forest per fold,
and timed each rule with a fresh prediction cache on one pinned CPU.
Each timed rule ran first, second and third once per candidate set. Both searches returned the
exhaustive vote's nominee on the 46 original sets and on 138 sets from
three draws of local feature exchanges around ANOVA-ranked anchors on the same folds. Exhaustive voting used 760
and 1,334 masks on the two original panels. Best-first search cut the
depth-first count by a further 10\% and 18\%. Nomination time fell by 9\%
and 17\%, respectively. A single point was faster on both panels
(Table~\ref{tab:minimum-cost-outcome}).
Best-first used fewer masked predictions than depth-first on 41 of 184 lists and tied
on 143. Of the 173 nonempty lists, 90 met the prefix minimum in Section~\ref{sec:certificate}.
On the other 83 the first candidate won and every rival scored zero. Those lists took
718 queries against a shortest total of 602. The clocks cover nomination, excluding ranking, model fitting,
the guard and the final refit.

\paragraph{Choice of subset.}
The single-point rule maximizes the smallest normalized tolerance margin
across metrics on the original panel with no features restored.
Selecting this score across all candidates increased
\(U\) over the minimum vote by 3.44 percentage points on the ten tasks
(\([-4.67,15.00]\)) and 5.49 on the six wider tasks (\([0.00,11.05]\)).
On the first panel its breach rate rose by 6.67 points. The first comparison
was sensitive to ties. Choosing uniformly among the highest-vote candidates
gave \(5.42\%\) effective compression, 2.53 points above the fixed order.

To isolate the contribution of voting with a common tie rule, we used
XGBoost for masking and refitting and retained the better half of each
fold's candidate sets by Gain rank sum. We compared the highest point
score on that shortlist with the highest point score among its
maximum-vote candidates. Voting changed only the Seeds nomination on
the ten tasks, improving \(U\) by 1.33 points (\([0,4.00]\)) with the
same breach rate. It changed none of the sixteen wider-fold nominations
and increased their total mask count from 71 to 168. This retrospective
comparison reused the measured outcomes. On these shortlists the vote
rarely altered the decision that received the refit.
